\documentclass[12pt,oneside]{article} % Uma Coluna e lingua portuguesa
%\usepackage[T1]{fontenc}        % Permite digitar os acentos de forma normal
\usepackage[utf8]{inputenc}
%\usepackage[english]{babel}
\usepackage[portuges,brazil]{babel}
%\usepackage[latin1]{inputenc}
\usepackage[dvips]{graphicx}    % Permite Gráficos
%\usepackage{times}    % Fonte Times
\usepackage{fancyhdr}
\usepackage{array}
\usepackage{multicol}
\usepackage[colorlinks=true,linkcolor=blue,urlcolor=blue]{hyperref}
\usepackage{nomencl}    % glossario
\usepackage{amssymb}
\usepackage{amsmath}
\usepackage[compact]{titlesec}
\usepackage{wrapfig}

%=======================================================================

% Hifenização das palavras desconhecidas pelo LaTeX
%\hyphenation{}
\paperheight    297mm
\paperwidth     210mm
\voffset         -15mm
\headheight      15pt %% tamanho de letra
\headsep         5mm  %% para o início do texto
\oddsidemargin  -3.0mm
\evensidemargin -3.0mm
\textwidth      167.0mm
\topmargin      005.0mm
\textheight     240.0mm
\footskip       10.0mm

\title{CI097 - Prova 1}

\author{Maratona de Programação}
\date{5 de Setembro de 2016}
\usepackage{indentfirst}
\usepackage{subfig}

\parindent=0pt
\setlength{\parskip}{7pt plus 1pt minus 2pt}
\titlespacing{\section}{0pt}{*0}{*0}
\titlespacing{\subsection}{0pt}{*0}{*0}
\titlespacing{\subsubsection}{0pt}{*0}{*0}

\begin{document}

\begin{center}
\textbf{\Huge CI097 - Prova 1}\\
\vspace{0.2cm}
\textit{5 de Setembro de 2016} \\
\vspace{1.0cm}
\textbf{Servidor BOCA:} \\
\texttt{\large http://maratona.c3sl.ufpr.br/} \\
\vspace{1.0cm}
\end{center}

\pagestyle{fancy}
\renewcommand{\footrulewidth}{0.7pt}
\renewcommand{\headrulewidth}{0.7pt}
\lhead{CI097 - Prova 1}
%\chead{Maratonas de Programação}
\rhead{5 de Setembro de 2016}
\cfoot{\thepage}


% Espaco para o create-zips.sh nao achar
  \section*{Instruções Importantes}

\begin{itemize}

    \item A prova é composta por 5 problemas que valem 20 pontos cada, e um
    problema bônus que também vale 20 pontos.

    \item O problema F é um problema bônus, usado na seletiva mas que não foi resolvido.
    Este problema só deve ser considerado caso todos os demais forem resolvidos.

    \item Se sua solução para um problema for aceita, a nota integral
    do problema é garantida.

    \item Se sua solução não for aceita, o professor irá avaliar sua última (mas
            \textbf{apenas} a última) submissão para o problema e atribuir a ela uma
    nota parcial, considerando a corretude do algoritmo utilizado, a
    complexidade computacional do algoritmo, o uso correto de funções e estruturas de dados
    disponíveis utilizadas, e eventuais detalhes da solução.

    \item Em cada problema, cada arquivo de entrada contém \textit{apenas um} caso de
    teste. Sua solução será executada com vários arquivos de entrada
    separadamente.

    \item Sua solução será compilada ou executada com a seguinte linha de comando:
    \begin{itemize}
        \item C: \verb|gcc -static -O2 -lm|
        \item C++: \verb|g++ -static -O2 -lm|
        \item C++11: \verb|g++ -std=c++11 -static -O2 -lm|
        \item Java: \verb|javac|
        \item Pascal: \verb|fpc -Xt -XS -O2|
        \item Python: \verb|python3|
    \end{itemize}

    \item O tempo limite para a execução de sua solução é de 1 segundo para os
    problemas A, B, C, D e E, e de 2 segundos para o problema F.

    \item Todas as linhas, tanto na entrada quanto na saída, terminam com o
    caractere de fim-de-linha ($\backslash n$), mesmo quando houver apenas uma única
    linha no arquivo.

    \item Para submissões em \textbf{JAVA}, a classe deverá ter o mesmo nome que
    o \textit{basename} do problema (leia a linha entre o título e o texto do
    problema).

\end{itemize}

\newpage

\clearpage
\section*{A: Transfira para o Professor} %tle=1
\vspace{-0.52cm}
\noindent \begin{verbatim}Arquivo: transfira.[c|cpp|java|pas|py]\end{verbatim}
\textit{Analógimôn Go} é um jogo bastante popular. Em sua jornada, o jogador percorre diversas
cidades capturando pequenos monstrinhos virtuais, chamados \textit{analógimôns}.

Você é um experiente jogador e já capturou $N$ analógimôns, numerados de $1$ a
$N$. Você capturou tantos monstrinhos que já está difícil levar todos com você
em sua jornada. Por isso, você pode se livrar de alguns de seus
monstrinhos transferindo-os para o Professor.

Ao transferir o analógimôn $i$ (para $1 \leq i \leq N$) para o Professor, você ganha
$D_i$ doces do Professor em troca do monstrinho. Como os doces são itens muito
importantes no jogo, você
quer transferir quais e quantos analógimôns forem necessários para ter a maior
quantidade possível de doces!

Entretanto, o analógimôn $i$ (para $1 \leq i \leq N$) pesa $P_i$ kg, e, devido a uma limitação de
espaço no laboratório do Professor, ele não pode receber analógimôns cuja soma
total dos pesos é maior que $K$ kg.

Sua tarefa é determinar a quantidade máxima de doces que você pode obter
transferindo seus monstrinhos, respeitando a limitação de espaço do laboratório do Professor.


\subsection*{Entrada}

A primeira linha contém os inteiros $N$ e $K$
($1 \leq N \leq 100, 1 \leq K \leq 10^4$), o número de analógimôns que você
capturou e a
capacidade do laboratório do Professor, em kg, respectivamente.
A segunda linha contém $N$ inteiros $D_1$, ..., $D_N$ ($1 \leq D_i \leq 10^4$), indicando quantos doces
você ganhará pela transferência de cada analógimôn. A terceira linha contém $N$
inteiros $P_1$, ..., $P_N$ ($1 \leq P_i \leq 10^4$) indicando o peso de cada analógimôn, em kg.

\subsection*{Saída}

Imprima uma única linha contendo a quantidade máxima de doces que você pode
obter.

\begin{table}[!h]
\centering
\begin{tabular}{|l|l|}
\hline
\begin{minipage}[t]{3in}
\textbf{Exemplo de entrada}
\begin{verbatim}
4 52
1 8 14 22
4 12 20 30
\end{verbatim}
\vspace{1mm}
\end{minipage}
&

\begin{minipage}[t]{3in}
\textbf{Exemplo de saída}
\begin{verbatim}
36
\end{verbatim}
\vspace{1mm}
\end{minipage} \\
\hline
\end{tabular}
\end{table}

\begin{table}[!h]
\centering
\begin{tabular}{|l|l|}
\hline
\begin{minipage}[t]{3in}
\textbf{Exemplo de entrada}
\begin{verbatim}
3 2
9 5 2
12 8 42
\end{verbatim}
\vspace{1mm}
\end{minipage}
&

\begin{minipage}[t]{3in}
\textbf{Exemplo de saída}
\begin{verbatim}
0
\end{verbatim}
\vspace{1mm}
\end{minipage} \\
\hline
\end{tabular}
\end{table}


\clearpage
\section*{B: Impeachment do Líder} %tle=1
\vspace{-0.52cm}
\noindent \begin{verbatim}Arquivo: impeachment.[c|cpp|java|pas|py]\end{verbatim}
\textit{Analógimôn Go} é um jogo bastante popular.
Os jogadores de \textit{Analógimôn Go} são divididos em três grandes
times: Time \textit{Valor}, Time \textit{Instinto} e Time \textit{Místico}, que
são liderados pelos seus líderes Kandera, Esparky e Blanque, respectivamente.
Naturalmente, você faz parte de um desses times!

O líder do seu time está sendo acusado de infringir as regras do jogo por
gerenciar incorretamente os doces recebidos do Professor que são destinados ao
time. Isto criou uma grande polêmica dentro da equipe: alguns jogadores defendem
que o líder realmente agiu incorretamente e deve sofrer um \textit{impeachment}
e ser afastado de seu cargo, enquanto outros defendem que ele não infringiu as
regras, que a acusação é inverídica e que ele deve continuar no cargo.

Para resolver a situação, uma votação será realizada entre todos os $N$
jogadores do seu time. Cada jogador deverá votar se o \textit{impeachment} deve
ou não ocorrer. Se o número de votos favoráveis ao \textit{impeachment} foi
maior ou igual a 2/3 (dois terços) do total de jogadores, o líder será afastado.
Caso contrário, a acusação é arquivada e ele continuará no cargo.

Dados os votos de todos os jogadores, determine o resultado da votação.

\subsection*{Entrada}

A primeira linha contém o inteiro $N$ ($1 \leq N \leq 10^5$), o número de
jogadores em seu time.
A próxima linha contém $N$ inteiros $v_1$, ..., $v_N$ ($v_i = 0$ ou $1$),
indicando os votos dos jogadores. O valor $1$ indica um voto favorável ao
\textit{impeachment}, enquanto o valor $0$ indica um voto contrário ao mesmo.

\subsection*{Saída}

Imprima uma linha contendo a palavra \verb|impeachment| se o líder deve
ser afastado de seu cargo, ou \verb|acusacao arquivada| caso contrário.

\begin{table}[!h]
\centering
\begin{tabular}{|l|l|}
\hline
\begin{minipage}[t]{3in}
\textbf{Exemplo de entrada}
\begin{verbatim}
6
1 0 1 1 0 1
\end{verbatim}
\vspace{1mm}
\end{minipage}
&

\begin{minipage}[t]{3in}
\textbf{Exemplo de saída}
\begin{verbatim}
impeachment
\end{verbatim}
\vspace{1mm}
\end{minipage} \\
\hline
\end{tabular}
\end{table}

\begin{table}[!h]
\centering
\begin{tabular}{|l|l|}
\hline
\begin{minipage}[t]{3in}
\textbf{Exemplo de entrada}
\begin{verbatim}
5
0 1 1 1 0
\end{verbatim}
\vspace{1mm}
\end{minipage}
&

\begin{minipage}[t]{3in}
\textbf{Exemplo de saída}
\begin{verbatim}
acusacao arquivada
\end{verbatim}
\vspace{1mm}
\end{minipage} \\
\hline
\end{tabular}
\end{table}


\clearpage
\section*{C: Ginásio} %tle=1
\vspace{-0.52cm}
\noindent \begin{verbatim}Arquivo: ginasio.[c|cpp|java|pas|py]\end{verbatim}
\textit{Analógimôn Go} é um jogo bastante popular. Em sua jornada, o jogador percorre diversas
cidades capturando pequenos monstrinhos virtuais, chamados \textit{analógimôns}.
As cidades contém localidades especiais chamadas de
\textit{ginásios}. Ao chegar a um ginásio, um jogador pode \textit{tentar}
colocar um de seus analógimôns dentro dele.

Cada anlógimôn tem dois inteiros associados a ele: seu \textit{Poder de Combate}
(PC) e seu \textit{Número de Ataques} (NA). Além disso, um ginásio tem associado
a ele um \textit{Intervalo de Poder} (IP). Ao tentar colocar um
analógimôn em um ginásio, o jogo verifica quantos são os analógimôns já
presentes no ginásio cuja diferença do seu PC para o PC do analógimôn sendo
colocado é de no máximo IP. Se esta quantidade for menor ou igual
ao NA do analógimôn sendo colocado, o monstrinho é inserido no ginásio com
sucesso. Caso contrário, ele não é colocado no ginásio.
Em ambos os casos, os analógimôns que já estavam no ginásio
continuam no ginásio.
Como exemplo, considere um ginásio com IP=$3$ com analógimôns de PC iguais a
$5$, $8$,
$13$ e $20$. Se um jogador tenta colocar um analógimôn de PC=$10$ e NA=$4$, o jogo
contará quantos analógimôns há no ginásio com PC entre $10-3 = 7$ e $10+3 = 13$,
inclusive. Como há dois analógimôns neste caso, o monstrinho é colocado com
sucesso no ginásio, pois $2 \leq 4$. O ginásio passa a conter analógimôns de PC
iguais a $5$, $8$, $10$, $13$ e $20$.

Dadas as informações sobre um ginásio e as tentativas de colocar analógimôns
dentro dele, determine quantos analógimôns ficarão no ginásio após todas as
tentativas. Considere que o ginásio inicialmente não contém nenhum analógimôn.

\subsection*{Entrada}

A primeira linha contém os inteiros $IP$ e $M$ ($1 \leq IP, M \leq 10^5$), o
IP do ginásio e o número de tentativas, respectivamente.
As próximas $M$ linhas descrevem as tentativas de colocar um analógimôn no
ginásio, na ordem em que são feitas. Cada linha contém dois inteiros $PC$ e $NA$
($1 \leq PC, NA \leq 10^5$), indicando o PC e o NA
do analógimôn, respectivamente. O Poder de Combate de todos analógimôns são
distintos.

\subsection*{Saída}

Imprima uma linha com um inteiro indicando quantos analógimôns
ficarão no ginásio.

\begin{table}[!h]
\centering
\begin{tabular}{|l|l|}
\hline
\begin{minipage}[t]{3in}
\textbf{Exemplo de entrada}
\begin{verbatim}
3 7
5 2
13 1
8 1
20 5
6 1
11 1
10 4
\end{verbatim}
\vspace{1mm}
\end{minipage}
&

\begin{minipage}[t]{3in}
\textbf{Exemplo de saída}
\begin{verbatim}
5
\end{verbatim}
\vspace{1mm}
\end{minipage} \\
\hline
\end{tabular}
\end{table}


\clearpage
\section*{D: Tipos de Espécies} %tle=1
\vspace{-0.52cm}
\noindent \begin{verbatim}Arquivo: tipos.[c|cpp|java|pas|py]\end{verbatim}
\textit{Analógimôn Go} é um jogo bastante popular. Em sua jornada, o jogador percorre diversas
cidades capturando pequenos monstrinhos virtuais, chamados \textit{analógimôns}.
Existem várias espécies de analógimôns. Cada espécie é de
(exatamente) um \textit{tipo}, como \textit{fogo}, \textit{água}, \textit{elétrico}, etc.
Algumas espécies podem ser do mesmo tipo, enquanto outras podem se tipos
diferentes.

No manual oficial do jogo consta que algumas espécies são do
mesmo tipo. Entretanto, o manual pode não apresentar esta informação para todos
os pares de espécies que são do mesmo tipo. Por exemplo, se o manual indica
que uma espécie $a$ é do mesmo tipo que uma espécie $b$, e que uma espécie $b$ é
do mesmo tipo que uma espécie $c$, então as espécies $a$ e $c$ certamente são do mesmo
tipo, embora esta informação pode não constar no manual.

Você capturou um analógimôn de uma certa espécie. Sua tarefa é determinar o
menor número possível de espécies que certamente são do mesmo tipo da espécie
do seu analógimôn, de acordo com as informações contidas no manual.

\subsection*{Entrada}

A primeira linha contém dois inteiros $N$ e $M$ ($1 \leq N \leq 1000, 0 \leq M
\leq \frac{N \times (N-1)}{2})$, o número de espécies de analógimôns e o
número de informações presentes no manual, respectivamente. As espécies são
numeradas de $1$ a $N$.
Cada uma das próximas $M$ linhas contém uma informação presente no manual. Cada
linha contém dois inteiros $a$ e $b$ ($1 \leq a, b \leq N, a \neq b$), indicando
que as espécies $a$ e $b$ são do mesmo tipo.
A última linha contém um inteiro $E$ ($1 \leq E \leq N$), indicando a espécie de
seu analógimôn.

\subsection*{Saída}

Imprima uma linha com um inteiro indicando a menor quantidade de espécies de
analógimôns que certamente são do mesmo tipo da espécie do seu analógimôn, de
acordo com o manual. Note que a espécie do seu analógimôn também deve
ser contada.

\begin{table}[!h]
\centering
\begin{tabular}{|l|l|}
\hline
\begin{minipage}[t]{3in}
\textbf{Exemplo de entrada}
\begin{verbatim}
5 3
1 3
3 5
2 4
1
\end{verbatim}
\vspace{1mm}
\end{minipage}
&

\begin{minipage}[t]{3in}
\textbf{Exemplo de saída}
\begin{verbatim}
3
\end{verbatim}
\vspace{1mm}
\end{minipage} \\
\hline
\end{tabular}
\end{table}

\begin{table}[!h]
\centering
\begin{tabular}{|l|l|}
\hline
\begin{minipage}[t]{3in}
\textbf{Exemplo de entrada}
\begin{verbatim}
3 1
1 2
3
\end{verbatim}
\vspace{1mm}
\end{minipage}
&

\begin{minipage}[t]{3in}
\textbf{Exemplo de saída}
\begin{verbatim}
1
\end{verbatim}
\vspace{1mm}
\end{minipage} \\
\hline
\end{tabular}
\end{table}


\clearpage
\section*{E: O Último Analógimôn} %tle=1
\vspace{-0.52cm}
\noindent \begin{verbatim}Arquivo: ultimo.[c|cpp|java|pas|py]\end{verbatim}
\textit{Analógimôn Go} é um jogo bastante popular. Em sua jornada, o jogador percorre diversas
cidades capturando pequenos monstrinhos virtuais, chamados \textit{analógimôns}.
Você acabou de chegar em uma cidade que contém o último analógimôn que falta
para sua coleção!

A cidade pode ser descrita como um \textit{grid} de $N$ linhas e $M$ colunas.
Você está em uma dada posição da cidade, enquanto o último analógimôn está em
outra posição da mesma cidade. A cada segundo, você pode se mover (exatamente) uma
posição ao norte, ao sul, a leste ou a oeste. Considerando que o analógimôn não
se move, sua tarefa é determinar o menor tempo necessário para
ir até a posição do monstrinho.
A figura descreve o exemplo da
entrada, e apresenta um caminho percorrido em 5 segundos. Outros caminhos
percorridos no mesmo tempo são possíveis, mas não há outro caminho que pode ser
percorrido em um tempo menor.

\begin{center}
\includegraphics[scale=0.4]{ultimo/ultimo.pdf}
\end{center}

\subsection*{Entrada}

A primeira linha contém dois inteiros $N$ e $M$ ($2 \leq N, M \leq 100$), o
número de linhas e de colunas na cidade, respectivamente. As próximas $N$ linhas
contém $M$ inteiros cada, descrevendo a cidade. O inteiro $0$ indica uma posição
em branco; o inteiro $1$ indica a sua posição na cidade; o inteiro $2$ indica a
posição do analógimôn na cidade.
É garantido que haverá exatamente um inteiro
$1$ e exatamente um inteiro $2$ na descrição da cidade, e que os demais inteiros
serão iguais a $0$.

\subsection*{Saída}

Imprima uma linha contendo o menor tempo necessário para ir até o
monstrinho, em segundos.

\begin{table}[!h]
\centering
\begin{tabular}{|l|l|}
\hline
\begin{minipage}[t]{3in}
\textbf{Exemplo de entrada}
\begin{verbatim}
4 4
2 0 0 0
0 0 0 0
0 0 0 0
0 0 1 0
\end{verbatim}
\vspace{1mm}
\end{minipage}
&

\begin{minipage}[t]{3in}
\textbf{Exemplo de saída}
\begin{verbatim}
5
\end{verbatim}
\vspace{1mm}
\end{minipage} \\
\hline
\end{tabular}
\end{table}


%\clearpage
%\section*{F: Analógimôn Go} %tle=2
%\vspace{-0.52cm}
%\noindent \begin{verbatim}Arquivo: analogimongo.[c|cpp|java|pas|py]\end{verbatim}
%\input{analogimongo/analogimongo.tex}

\end{document}
